\chapter{Realisierung}

\section{Website Fingerprinting Versuchsaufbau Stand der Technik}

Als Voraussetzung wird ein Webserver benötigt welche eine Menge an definierten Websiten bereitstellt.
In dem Paper von Jansen wurde aufgezeigt dass dies mit der zimply Python Library erreicht werden kann.
Diese Bibliothek stellt einen Webserver für ZIM Dateien zur Verfügung.
\cite[S. 4]{jansen2023data} 
\footnote{\url{https://github.com/kimbauters/ZIMply}}


Jansen und Wails haben gezeigt, dass \texttt{wget2} als Surrogat für den Tor Browser verwendet werden kann.
Sowohl im Live-Tor-Netzwerk als auch in einer Shadow-Simulation wiesen die erzeugten Tor-Zellsequenzen hinreichende Ähnlichkeit für Website-Fingerprinting-Analysen auf. \cite[S.~6]{jansen2023data}
